\documentclass[12pt]{article}
\usepackage[utf8]{inputenc}
\usepackage[T1]{fontenc}
\usepackage{hyperref}
\title{Sause Language Resources for K-12 Curriculum Development}
\author{Compiled by Claude Sonnet 5.5 (claude-sonnet-5-5) for Andr\'as Kornai}
\date{}
\begin{document}
\maketitle

\section{Introduction}
Sause (ISO 639-3 \texttt{sao}; Glottocode saus1247) is a Papuan language of Indonesia with about 250 speakers. It is spoken southwest of Sentani in the villages of Ures, Mubararon, Sause-Bokoko, Witti-Yadow, Lidya and Puaral. Usher groups it with Kapauri as a branch of Kwerbic, within a Tor-Kwerba proposal, but its affiliation is debated. The searches found only encyclopedic and linguistic overview pages.

\section{Primary Educational Resources}
No primers, textbooks or school programmes were found.

\section{Linguistic and Reference Materials}
\begin{itemize}
\item Wikipedia on Sause: \url{https://en.wikipedia.org/wiki/Sause_language}
\item Wikipedia on the Kapauri--Sause languages: \url{https://en.wikipedia.org/wiki/Kapauri%E2%80%93Sause_languages}
\item Wikipedia on the Kwerbic languages: \url{https://en.wikipedia.org/wiki/Kwerbic_languages}
\item New Guinea World, Kapauri-Sause family page: \url{https://newguineaworld.linguistik.uzh.ch/families/northwest-new-guinea/foja-range/west-foja-range/kapauri-sause/start}
\end{itemize}

\section{Religious and Community Resources}
None were found.

\section{Digital and Audio Resources}
None were found.

\section{Curriculum Development Context}
No teaching or documentation materials were found. A curriculum would start from new documentation and community work. Indonesian is the national language of instruction.

\end{document}